
\subsubsection{3.1 Stratified Training for Retrieval-Quality Classification}

Our core methodological innovation is stratified sampling that enforces divergence from educational quality. We extract positive examples (BEIR qrel scores $\geq 1)$ and negative examples (qrel scores = 0) from BEIR Natural Questions. For each example, we compute educational quality (GPT-2 perplexity, lower is higher quality) and BEIR quality (relevance score). We identify divergent examples---documents with high perplexity (low educational quality, above median) but high BEIR relevance (above median)---and oversample them $3\times$ during training.

This stratification forces classifiers to learn signals orthogonal to fluency. Without stratification, classifiers would likely rediscover perplexity. By oversampling low-perplexity, high-BEIR documents, we decorrelate fluency from relevance. If classifiers achieve validation accuracy >70\% despite this decorrelation, they have learned signals beyond perplexity.

We use FastText \citep{Joulin2017} for computational efficiency at corpus scale: embedding dimension 100, learning rate 0.1, 25 epochs, word n-grams 2. Training completes in under 5 seconds on CPU.

\subsubsection{3.2 Entity Density Measurement}

We operationalize factual density through named entity recognition, hypothesizing that retrieval-optimal documents contain more entities per unit text. We use spaCy's en\_core\_web\_sm model to extract factual entity types (PERSON, ORGANIZATION, GPE, DATE, MONEY). Entity density is computed as entities per 100 tokens, normalized for document length.

For mechanism testing, we measure entity density in retrieval-classifier-selected corpus (5,000 documents) and perplexity-baseline corpus (5,000 documents, matched size). The ratio $\rho$ = density\_retrieval / density\_perplexity operationalizes factual informativeness. A ratio $\geq 1.15$ would indicate the classifier learned entity-based features; a ratio <1.15 falsifies the density-learning hypothesis.

\subsubsection{3.3 Query Splitting for Differential Evaluation}

To test whether entity density mechanistically drives retrieval gains, we split queries by BM25 performance. We run BM25 (Okapi, k1=1.5, b=0.75) and classify queries as lexical (relevant document in BM25 top-10) or semantic (BM25 fails). If entity density drives quality, high-density corpora should preferentially help semantic queries where entities provide alternative lexical pathways.

\subsubsection{3.4 Experimental Design}

We designed three experiments:

\textbf{H-E1 (Existence):} Does retrieval-quality filtering improve Recall@10 over perplexity filtering? Gate threshold: $\geq +0.03pp$ improvement. Tests whether retrieval-specific filtering is feasible at proof-of-concept scale.

\textbf{H-M1 (Mechanism: Density):} Do classifier-selected documents exhibit higher entity density? Gate threshold: ratio $\geq 1.15$ (15\% increase). Tests whether stratified training teaches classifiers to prioritize factual density.

\textbf{H-M2 (Mechanism: Selectivity):} Do high-density documents preferentially improve semantic queries? Gate threshold: $\Delta Recall_semantic \geq 0.04pp$ and $\Delta Recall_lexical \leq 0.01pp$ (3pp differential). Tests whether improvements derive from entity-based semantic quality versus lexical coverage.

This decomposition separates existence claims from mechanistic claims. Falsifying H-M1 or H-M2 while validating H-E1 would indicate retrieval-quality signals exist but operate through mechanisms other than entity density.
